\section{Proyek Mini-Kompilator dan Strategi Belajar}
\label{sec:ch01-proyek}

Sepanjang semester, mahasiswa membangun satu proyek berkelanjutan:
mini-kompilator untuk bahasa \textbf{Subset-C} (SC). Proyek ini tumbuh
bersama bab. Pada Bab VII mahasiswa menyelesaikan \eng{lexer}; pada
Bab X--XIII menyelesaikan \eng{parser} dan AST; pada Bab XIV--XVI
menambahkan analisis semantik, tabel simbol, pemeriksaan tipe, dan
validator.

\begin{table}[H]
\centering
\caption{Milestone proyek dan bab pendukungnya}
\label{tab:ch01-milestone}
\begin{tabularx}{\textwidth}{@{}llX@{}}
\toprule
\textbf{Milestone} & \textbf{Bab} & \textbf{Keluaran} \\
\midrule
M0 & III & Struktur repo, spesifikasi bahasa v0.1, program contoh \\
M1 & VII & Lexer tulisan tangan dan lexer Flex, 30 tes \\
M2 & X & Recursive descent parser, AST minimal, tes \\
M3 & XIII & Grammar Bison, AST lengkap, integrasi build \\
M4 & XV & Tabel simbol bersarang dan pemeriksaan tipe \\
M5 & XVI & Validator, diagnostik, dan proyek akhir \\
\bottomrule
\end{tabularx}
\end{table}

Strategi belajar yang disarankan adalah \emph{baca, kerjakan, jelaskan}.
Baca materi sebelum kuliah, kerjakan praktikum dan latihan secara
bertahap, lalu jelaskan hasilnya kepada teman atau tuliskan dalam
jurnal belajar. Pola ini melatih pemahaman konsep sekaligus kemampuan
komunikasi teknis.

\begin{tipsbox}
Jangan menunda praktikum sampai akhir semester. Setiap milestone hanya
memerlukan tambahan kecil di atas milestone sebelumnya. Jika proyek
dikerjakan konsisten, proyek akhir pada Bab XVI akan terasa sebagai
penyelesaian alami, bukan pekerjaan mendadak.
\end{tipsbox}
